\section{Related Work}

A growing body of financial natural language processing benchmarks has
emerged, but nearly all of them evaluate \emph{comprehension} over
existing reports rather than \emph{construction} of new accounting
artifacts.
FinQA~\citep{chen2021finqa} studies numerical reasoning over
earnings reports;
TAT-QA~\citep{zhu2021tatqa} extends the setting to joint table--text
reasoning; and
DocFinQA~\citep{reddy2024docfinqa} situates similar questions within
long-form SEC filings.
These benchmarks require models to read a financial document and answer
questions about it, but not to construct a balance sheet from raw
transactions.

Broader evaluation suites such as
FinanceBench~\citep{islam2023financebench} cover wider portions of
financial language understanding and decision support, while
FinanceMATH~\citep{zhao2024financemath} emphasizes exact quantitative
reasoning in finance settings.
Even MMLU~\citep{hendrycks2021mmlu}, which includes
professional-accounting questions, measures declarative knowledge in a
multiple-choice format rather than procedural bookkeeping.

The closest adjacent line of work studies richer financial analysis
over existing statements or filings rather than bookkeeping from raw
events.
FinAuditing
~\citep{wang2025finauditing} studies taxonomy-structured auditing over
multi-document financial disclosures.
Wang et al.~\citep{wang2025automating} also study automated financial
statement auditing with a benchmark that combines financial tables and
synthesized transaction data.
These tasks are closer to professional finance workflows than QA alone,
but they primarily begin from prepared financial artifacts and evaluate
analysis, diagnosis, or revision rather than direct construction of the
final statement from journal entries and adjustments.
More general multi-step arithmetic benchmarks such as
GSM8K~\citep{cobbe2021gsm8k} test quantitative reasoning, but they do
not impose the domain-specific conservation constraints, section
structure, and closing-entry behavior required in accounting.

Our benchmark therefore occupies a distinct point in the literature.
It evaluates whether a model can transform raw accounting events into a
complete, internally consistent balance sheet with exact ground truth.
That setting matters because exact invariants reveal failures that are
easy to hide in answer-only evaluation: a model may recover the right
accounts, sections, and financial vocabulary while still breaking the
accounting equation or mishandling equity closure.
